% ============================================================================
\section{初始化：\texttt{syscall\_init()}}
\label{sec:syscall-init}
% ============================================================================

\texttt{syscall\_init()} 做三件事：清空 table、注册 handler、调用后端 init。

\codefile{src/kernel/core/syscall.c}
\begin{lstlisting}[style=cstyle]
void syscall_init(void) {
    /* 1. 清空 syscall table */
    for (unsigned i = 0; i < SYSCALL_MAX; i++) {
        syscall_table[i] = NULL;
    }

    /* 2. 注册内置 syscall handler */
    syscall_register(SYS_EXIT,  sys_exit);
    syscall_register(SYS_WRITE, sys_write);
    syscall_register(SYS_BRK,   sys_brk);

    printk("[syscall] table: %u handlers registered "
           "(exit=%u, write=%u, brk=%u)\n",
           3, SYS_EXIT, SYS_WRITE, SYS_BRK);

    /* 3. 初始化后端（安装 IDT gate 或配置 MSR） */
    if (g_syscall_ops) {
        printk("[syscall] backend: %s\n", g_syscall_ops->name);
        g_syscall_ops->init();
    } else {
        printk("[syscall] no backend registered "
               "(dispatch only)\n");
    }

    g_syscall_ready = 1;
    printk("[syscall] init ok\n");
}
\end{lstlisting}

% ----------------------------------------------------------------------------
\subsection{kmain.c 集成}
% ----------------------------------------------------------------------------

\texttt{syscall\_init()} 在 VMA 初始化之后调用。
这个顺序有讲究：

\codefile{src/kernel/core/kmain.c}
\begin{lstlisting}[style=cstyle]
        vma_init();

        if (vma_is_ready()) {
            vma_add(KERNEL_VIRT_OFFSET, vmm_direct_map_end(),
                    VMA_READ | VMA_WRITE | VMA_EXEC,
                    "direct-map");
            vma_add(KHEAP_START,
                    KHEAP_START + KCONFIG_HEAP_INITIAL_PAGES
                                 * VMM_PAGE_SIZE,
                    VMA_READ | VMA_WRITE, "kheap");
        }

        /*
         * Stage D-2: 系统调用子系统初始化
         *
         * 在 VMA 之后，因为 sys_write 需要通过 VMA
         * 验证用户态指针的合法性。
         */
        syscall_init();
\end{lstlisting}

\begin{whybox}
为什么 \texttt{syscall\_init} 要在 VMA 之后？

\texttt{sys\_write} 需要验证用户传入的 \texttt{buf} 指针不在内核空间。
将来更精确的检查会查询 VMA——确认 \texttt{[buf, buf+count)} 落在合法的用户 VMA 中。
如果 VMA 还没初始化，这些检查就无法执行。

另一方面，\texttt{syscall\_init} 不依赖后端——即使没有注册后端，
dispatch 层也能正常工作（只是没有用户态入口点）。
\end{whybox}

% ----------------------------------------------------------------------------
\subsection{\texttt{kconfig.h}：后端选择}
% ----------------------------------------------------------------------------

\codefile{src/include/kernel/kconfig.h}
\begin{lstlisting}[style=cstyle]
/*
 * 系统调用后端
 *
 *   0 = int0x80   -- 传统软中断（IDT[0x80] DPL=3）
 *   1 = sysenter  -- Intel 快速系统调用（MSR 配置）
 */
#define KCONFIG_SYSCALL_BACKEND  0
\end{lstlisting}

目前还没有后端实现，所以 \texttt{kmain.c} 中暂不调用
\texttt{syscall\_register\_backend()}。D-3 实现 \texttt{int 0x80} 后端后，
会在 \texttt{kmain.c} 中添加基于 \texttt{KCONFIG\_SYSCALL\_BACKEND} 的选择逻辑，
和 PMM/VMA 的后端选择完全一致。
